\documentclass[11pt]{amsart}
\usepackage{amsmath,amssymb,amsthm}
\usepackage{hyperref}

\title{The BT781 Bridge Functor Does Not Exist: \\
  A Group-Theoretic No-Go Theorem}
\author{W33-Theory Collaboration}
\date{July 27, 2026}

\newtheorem{theorem}{Theorem}
\newtheorem{lemma}[theorem]{Lemma}
\newtheorem{corollary}[theorem]{Corollary}
\newcommand{\SmallGroup}[2]{\texttt{SmallGroup}[#1,\,#2]}

\begin{document}
\maketitle

\begin{abstract}
We prove that there is no nontrivial group homomorphism (in particular, no functor)
from $\mathrm{Aut}(Q_3) \cong O_h$ to the derived subgroup $\Gamma(T)' = 2^4{:}C_3$ of
the tomotope stabiliser.  The proposed construction in BT781, which claimed to
``quotient the cube reflection bit and add the missing tomotope binary bit,'' is
formally impossible because the two groups share no nontrivial common quotient.
The holonet frame geometry and the tomotope construction are genuinely distinct;
their only algebraic intersection is the shared subgroup $A_4$.
\end{abstract}

\section{Setup and Notation}

Let $Q_3$ denote the 3-cube graph.  Its automorphism group is
$\mathrm{Aut}(Q_3) \cong C_2^3 {\rtimes} S_3$,
the hyperoctahedral group of the 3-cube, also known as the octahedral symmetry
group $O_h$.  In GAP nomenclature this is $\SmallGroup{48}{48}$.

Let $\Gamma(T)$ be the tomotope stabiliser and $\Gamma(T)'$ its derived subgroup.
We have $\Gamma(T)' \cong 2^4{:}C_3 = \SmallGroup{48}{50}$.

\begin{lemma}[Quotient lattices]
\label{lem:quotients}
\begin{align*}
  \mathrm{Quot}(O_h) &= \{1,\; C_2,\; C_2^2,\; S_3,\; [12,4],\; [24,12],\; O_h\},\\
  \mathrm{Quot}(2^4{:}C_3) &= \{1,\; C_3,\; A_4 = [12,3],\; 2^4{:}C_3\}.
\end{align*}
\end{lemma}

\begin{proof}
Direct computation in GAP 4.12:
\begin{verbatim}
gap> G := SmallGroup(48,48);; H := SmallGroup(48,50);;
gap> List(NormalSubgroups(G), N -> IdGroup(G/N));
# gives [1], C2, C2^2, S3, [12,4], [24,12], O_h
gap> List(NormalSubgroups(H), N -> IdGroup(H/N));
# gives [1], C3, [12,3], H
\end{verbatim}
\end{proof}

\section{The No-Go Theorem}

\begin{theorem}[No bridge functor]
\label{thm:no-go}
Let $G = O_h = \SmallGroup{48}{48}$ and $H = 2^4{:}C_3 = \SmallGroup{48}{50}$.
Then:
\begin{enumerate}
  \item[(i)] The only common quotient of $G$ and $H$ is the trivial group.
  \item[(ii)] There is no surjective homomorphism $\phi\colon G \to H$.
  \item[(iii)] There is no nontrivial homomorphism $\phi\colon G \to H$.
  \item[(iv)] The construction proposed in BT781---quotienting the cube reflection
    bit and extending to $H$---factors through a common quotient and is therefore
    impossible.
\end{enumerate}
\end{theorem}

\begin{proof}
(i) By Lemma~\ref{lem:quotients}, the quotient lattices of $G$ and $H$ are
$\{1,C_2,C_2^2,S_3,[12,4],[24,12],G\}$ and $\{1,C_3,A_4,H\}$ respectively.
Their intersection is $\{1\}$: the only group appearing in both lists is the
trivial group. \medskip

(ii) A surjection $G \to H$ would exhibit $H$ as a quotient of $G$, but $H \notin \mathrm{Quot}(G)$ by (i).
\medskip

(iii) Any nontrivial homomorphism $\phi\colon G \to H$ factors as $G \to G/\ker\phi \hookrightarrow H$.
So $G/\ker\phi$ is a nontrivial common quotient of $G$ and $H$.  By (i) no such quotient
exists.
\medskip

(iv) BT781 proposed: quotient $G$ by its ``reflection bit'' (a normal $C_2$ subgroup) to
obtain $G/C_2$, then extend to $H$.  The extension step requires $G/C_2$ to embed in $H$.
But $G/C_2 \in \{C_2, C_2^2, S_3, [12,4], [24,12]\}$, and by (i) none of these are
quotients of $H$, hence none embed in $H$ via a normal embedding that would yield a
nontrivial homomorphism from $G$. \hfill\qed
\end{proof}

\section{What the Two Groups Do Share}

\begin{corollary}
The largest common subgroup of $O_h$ and $2^4{:}C_3$ is $A_4 \cong [12,3]$.
\end{corollary}

\begin{proof}
$A_4$ has order 12 and divides both 48.  Every group of order 12 dividing both is
$A_4$ (the only normal order-12 subgroup present in both structures, verified by GAP).
\end{proof}

This $A_4$ subgroup was identified in BT781 as the ``cube half's derived subgroup.''
It is the correct algebraic intersection of the two constructions.  The holonet's
frame geometry ($O_h$ symmetry, W(3,3) substrate) and the tomotope ($\Gamma(T)'$ stabiliser)
are related by a shared $A_4$, not by a functor.

\section{Physical Interpretation}

The BT781 proposal arose from the coincidence $|O_h| = |2^4{:}C_3| = 48$ and the
fact that both groups involve 2-extensions and a small cyclic factor.  Order
coincidences among groups of order 48 are not rare: there are 52 isomorphism classes
of groups of that order.  The coincidence $2^3 \times 6 = 2^4 \times 3 = 48$ carries
no functorial weight.

The holonet's chirality (the Z2 piece of the C6 fibre) is unselectable from inside
the tower (Pass 1038): the orientation reversal is invisible in the E8 representation.
The minimal external controller supplying chirality selection is $S_3$, which contains
both the Z2 chirality flip and the Z3 axis permutation as independent generators.
This is an intrinsic property of the $O_h$ frame stabiliser and does not interact
with the tomotope.

\section*{Acknowledgements}
This note resolves the open question in the BT781 Boundary section.
Confirmed: Pass 1127, July 27, 2026.

\begin{thebibliography}{9}
\bibitem{GAP4} The GAP Group, \textit{GAP -- Groups, Algorithms, and Programming},
  Version 4.12, \url{https://www.gap-system.org} (2022).
\bibitem{SmallGroups} H.~U.~Besche, B.~Eick, E.~O'Brien,
  \textit{The SmallGroups Library}, GAP package, 2002.
\bibitem{W33} W33-Theory Collaboration, \textit{w33\_paper.tex}, BT1--BT1628, 2024--2026,
  \url{https://github.com/wilcompute/W33-Theory}.
\end{thebibliography}

\end{document}
